% This document was generated by an LLM.

\documentclass{essay}

\title{Transforming Between Differential and Integral Equations}
\author{}
\date{}

\begin{document}

\maketitle

\section{Introduction}

Differential equations and integral equations are usually taught as separate subjects, with separate techniques and separate literatures. Yet in a large and important class of problems they are not two subjects at all, but two representations of a single underlying object. A differential equation is a \emph{local} statement: it constrains the behavior of an unknown function in infinitesimal neighborhoods, and additional data---initial or boundary conditions---must be supplied separately to single out a solution. An integral equation is a \emph{global} statement: a single, self-contained relation in which the governing law and the auxiliary data are fused together.

The dictionary between the two representations is, at bottom, the Fundamental Theorem of Calculus. But the translation is far from a triviality, because the two forms have opposite analytic personalities. Differentiation is an unbounded, roughening operation, sensitive to perturbation; integration is a bounded, compact, smoothing operation with a rich spectral theory. Much of classical analysis---the Picard--Lindel\"of existence theorem, Sturm--Liouville eigenfunction expansions, the modern theory of weak solutions---consists precisely of translating a differential problem into integral form, doing the hard analysis where the operators are well behaved, and translating back.

This essay develops the correspondence carefully. Section~\ref{sec:volterra} carries out the basic transformation from an initial value problem to a Volterra integral equation in full detail, paying particular attention to two points that are often glossed over: how the interval of integration is chosen, and why the equivalence between the two formulations is genuinely two-directional. Section~\ref{sec:second-order} extends the construction to second-order equations, where the notion of a \emph{kernel} first appears. Section~\ref{sec:fredholm} treats boundary value problems, whose integral counterparts are Fredholm equations built from Green's functions. Section~\ref{sec:symmetry} examines why the correspondence is a genuine operator-theoretic duality rather than a formal rewrite, and Section~\ref{sec:equivalence} states precisely in what sense, and with what caveats, the two formulations are equivalent.

\section{From an Initial Value Problem to a Volterra Equation}
\label{sec:volterra}

\subsection{Setup}

Consider the initial value problem
\begin{equation}
y'(t) = f\bigl(t, y(t)\bigr), \qquad y(t_0) = y_0,
\label{eq:ivp}
\end{equation}
where $f$ is continuous on a suitable domain. The goal is to convert~\eqref{eq:ivp} into a statement of the form ``$y(t) = \text{something}$'' that contains no derivatives and no separate side condition.

\subsection{Integrating both sides: how the interval is chosen}

The first step is to integrate both sides of the differential equation---but over what interval? Here is the key point: \emph{the interval of integration is not a free choice; it is dictated by what we know and what we want.}

We know the value of $y$ at exactly one point, namely $t_0$. We want the value of $y$ at a generic point $t$. So we integrate \emph{from the point where information lives to the point where information is wanted}. The equation $y'(s) = f(s, y(s))$ holds for every $s$ in the interval of existence; integrating it over $s \in [t_0, t]$ gives
\begin{equation}
\int_{t_0}^{t} y'(s)\, ds \;=\; \int_{t_0}^{t} f\bigl(s, y(s)\bigr)\, ds.
\label{eq:integrated}
\end{equation}

Two remarks are in order before continuing.

\begin{remark}[The dummy variable]
The upper limit of integration is $t$, a variable. One cannot use the same symbol for the integration variable and for a limit of integration: an expression such as $\int_{t_0}^{t} f(t, y(t))\, dt$ is meaningless, since the variable of integration would collide with a fixed endpoint. The integration variable is therefore renamed to $s$. This is exactly the same discipline as writing $\ln x = \int_1^x \frac{1}{s}\, ds$ rather than the ill-formed $\int_1^x \frac{1}{x}\, dx$.
\end{remark}

\begin{remark}[Why integrating both sides is legitimate]
The differential equation asserts that two functions of $s$ are equal \emph{pointwise} on an interval. Equal functions have equal integrals over any subinterval, so~\eqref{eq:integrated} follows. The converse fails---equal integrals do not imply equal integrands---which is why the reverse direction of the equivalence (Section~\ref{subsec:converse}) requires the Fundamental Theorem of Calculus rather than mere ``un-integration.''
\end{remark}

\subsection{Absorbing the initial condition}

By the Fundamental Theorem of Calculus (Part 2, the ``net change'' form), the left-hand side of~\eqref{eq:integrated} evaluates to
\[
\int_{t_0}^{t} y'(s)\, ds = y(t) - y(t_0) = y(t) - y_0.
\]
Observe what has happened: \emph{the initial condition has entered the equation.} It was not appended as a side remark; it appeared as the boundary term of the FTC evaluation. This is precisely the mechanism by which the side condition is absorbed into the integral formulation. Solving for $y(t)$:
\begin{equation}
\boxed{\; y(t) = y_0 + \int_{t_0}^{t} f\bigl(s, y(s)\bigr)\, ds \;}
\label{eq:volterra}
\end{equation}

Equation~\eqref{eq:volterra} is a \emph{Volterra integral equation of the second kind}: the unknown $y$ appears both outside and inside the integral, and the upper limit of integration is variable. Read as a sentence, it says: \emph{the state at time $t$ equals the initial state plus the accumulated instantaneous change from $t_0$ up to $t$.} The variable upper limit is what makes the equation ``Volterra'': the value at time $t$ depends only on the history $[t_0, t]$, never on the future. Causality is visible in the limits of integration.

\subsection{What goes wrong with other choices of interval}

Suppose one instead integrated over some fixed interval $[a, b]$:
\[
\int_a^b y'(s)\, ds = \int_a^b f\bigl(s, y(s)\bigr)\, ds
\quad\Longrightarrow\quad
y(b) - y(a) = \int_a^b f\bigl(s, y(s)\bigr)\, ds.
\]
This is a \emph{true} statement but a useless one: it is a single scalar relation between two unknown numbers $y(a)$ and $y(b)$, and it does not determine the function $y$. To characterize the whole function one needs one equation \emph{for each value of $t$}---which is exactly what a variable upper limit provides. The transformation produces a one-parameter family of equations, one per $t$, packaged as a single functional equation.

Moreover, the lower limit must be $t_0$ specifically, because $t_0$ is the only point at which the boundary term produced by the FTC is a \emph{known} number rather than another unknown. Integrating from some other point $c$ would yield $y(t) = y(c) + \int_c^t f(s, y(s))\, ds$, which is correct but not closed, since $y(c)$ is unknown. \emph{The anchor point of the integral is the anchor point of the data.}

Nothing requires $t \ge t_0$, incidentally. For $t < t_0$ formula~\eqref{eq:volterra} still holds under the usual orientation convention $\int_{t_0}^{t} = -\int_{t}^{t_0}$; the equation describes the solution on both sides of the initial point.

\subsection{The equivalence is two-directional}
\label{subsec:converse}

We have shown that a solution of the initial value problem satisfies the integral equation. The converse deserves care, and it is where the Fundamental Theorem of Calculus earns its keep a second time.

\begin{proposition}
Let $f$ be continuous, and suppose $y$ is merely a \emph{continuous} function satisfying~\eqref{eq:volterra}. Then $y$ is continuously differentiable, satisfies $y'(t) = f(t, y(t))$, and satisfies $y(t_0) = y_0$.
\end{proposition}

\begin{proof}[Proof sketch]
Since $y$ and $f$ are continuous, the integrand $s \mapsto f(s, y(s))$ is continuous. By the Fundamental Theorem of Calculus (Part 1), the right-hand side of~\eqref{eq:volterra} is differentiable in $t$ with derivative $f(t, y(t))$; hence so is $y$, and $y'(t) = f(t, y(t))$. Setting $t = t_0$ in~\eqref{eq:volterra} makes the integral vanish (degenerate interval), recovering $y(t_0) = y_0$.
\end{proof}

Note the asymmetry in the hypotheses: only continuity of $y$ was \emph{assumed}, and differentiability was \emph{derived}. The integral equation self-improves the regularity of its solutions. This bootstrapping is the reason analysts prefer the integral form: one may search for solutions in the large, friendly space $C[t_0, t_0 + h]$ of continuous functions---complete under the supremum norm, and hence ideal for the contraction mapping argument of the Picard--Lindel\"of theorem---and differentiability comes out for free at the end. Indeed, Picard iteration,
\[
y_{n+1}(t) = y_0 + \int_{t_0}^{t} f\bigl(s, y_n(s)\bigr)\, ds,
\]
operates directly on the integral form; the integral operator is a contraction on a suitable space, whereas the differential operator is not even continuous.

\section{Second-Order Equations: Where the Kernel Comes From}
\label{sec:second-order}

Consider now the second-order initial value problem
\[
y''(t) = f\bigl(t, y(t)\bigr), \qquad y(t_0) = y_0, \quad y'(t_0) = v_0.
\]
Integrating once (the same logic as before, applied to $y''$) gives
\[
y'(u) = v_0 + \int_{t_0}^{u} f\bigl(s, y(s)\bigr)\, ds,
\]
and integrating again, over $u \in [t_0, t]$,
\[
y(t) = y_0 + v_0 (t - t_0) + \int_{t_0}^{t} \!\int_{t_0}^{u} f\bigl(s, y(s)\bigr)\, ds\, du.
\]
Both initial conditions have now been absorbed: \emph{each integration harvested one boundary term}. The double integral extends over the triangle $\{(u, s) : t_0 \le s \le u \le t\}$. Interchanging the order of integration (Fubini), the variable $u$ ranges over $[s, t]$ for fixed $s$, and since the integrand does not depend on $u$, the inner integral contributes only the length $(t - s)$:
\[
\int_{t_0}^{t} \!\int_{t_0}^{u} f\bigl(s, y(s)\bigr)\, ds\, du
= \int_{t_0}^{t} \left( \int_{s}^{t} du \right) f\bigl(s, y(s)\bigr)\, ds
= \int_{t_0}^{t} (t - s)\, f\bigl(s, y(s)\bigr)\, ds.
\]
Therefore
\begin{equation}
y(t) = y_0 + v_0 (t - t_0) + \int_{t_0}^{t} (t - s)\, f\bigl(s, y(s)\bigr)\, ds.
\label{eq:second-order-volterra}
\end{equation}

The factor $(t - s)$ is a \emph{kernel}: the fingerprint of double integration collapsed into a single integral. It is the $n = 2$ case of Cauchy's formula for repeated integration, in which $n$-fold integration produces the kernel $(t - s)^{n-1} / (n-1)!$. Physically, for $y'' = a(t)$, formula~\eqref{eq:second-order-volterra} says that the displacement at time $t$ due to an acceleration impulse at time $s$ is proportional to the elapsed time $t - s$: an impulse acts longer if it happened earlier. The kernel is the \emph{influence function} of the past on the present---the causal, initial-value counterpart of the Green's function that arises in the boundary-value setting, to which we now turn.

\section{Boundary Value Problems and Fredholm Equations}
\label{sec:fredholm}

Suppose that instead of initial conditions the auxiliary data are boundary conditions, as in
\[
y''(t) = f\bigl(t, y(t)\bigr), \qquad y(0) = y(1) = 0.
\]
The same program---integrate, absorb the data---now produces an integral equation with \emph{fixed} limits, a \emph{Fredholm equation}:
\[
y(t) = \int_0^1 G(t, s)\, f\bigl(s, y(s)\bigr)\, ds,
\]
where $G(t, s)$ is the \emph{Green's function} of the operator $d^2/dt^2$ subject to those boundary conditions. For this example,
\[
G(t, s) =
\begin{cases}
\,s\,(1 - t), & 0 \le s \le t, \\[2pt]
\,t\,(1 - s), & t \le s \le 1.
\end{cases}
\]
Once again the auxiliary data have been baked into the kernel. The structural correspondence is:
\begin{itemize}
\item \emph{Initial conditions} (data at one point; causal evolution) correspond to \emph{Volterra kernels} (variable limit; memory of the past only).
\item \emph{Boundary conditions} (data at two ends; a global constraint) correspond to \emph{Fredholm kernels} (fixed limits; every point interacts with every point).
\end{itemize}

\section{Why This Is a Genuine Symmetry, Not a Formal Rewrite}
\label{sec:symmetry}

The differential operator $L$ and the integral operator $K$ (integration against the Green's function) are literally inverses of one another: $K = L^{-1}$. Solving $Ly = f$ is the same as computing $y = Kf$. This is the same duality familiar from transfer functions: a differential equation in the time domain becomes an algebraic, multiplicative relation after transform, and the Green's function plays the role of the impulse response.

The transformation is useful because the two operators have opposite analytic personalities.

\emph{Differentiation is unbounded and roughening.} The operator $L$ amplifies high frequencies, is sensitive to noise, and is discontinuous as an operator between the natural function spaces; small perturbations of $y$ can make $Ly$ blow up.

\emph{Integration is bounded, compact, and smoothing.} The operator $K$ damps high frequencies and improves regularity. Compact operators enjoy a spectral theory---the Fredholm alternative; discrete eigenvalues accumulating only at zero---that closely mirrors finite-dimensional linear algebra.

Consequently the eigenvalue problem $Ly = \lambda y$ of Sturm--Liouville theory becomes $Ky = \mu y$ with $\mu = 1/\lambda$, and one may invoke the spectral theorem for compact self-adjoint operators. This is how one rigorously proves that a regular Sturm--Liouville problem possesses a complete orthonormal basis of eigenfunctions: the hard analysis lives on the integral side.

A concrete illustration: for $y'' + \lambda y = 0$ with $y(0) = y(\pi) = 0$, the eigenfunctions are $\sin(nt)$ with eigenvalues $\lambda_n = n^2$. Constructing the Green's function of $-\,d^2/dt^2$ on $[0, \pi]$ and writing the equivalent Fredholm equation, one finds that the same functions $\sin(nt)$ are eigenfunctions of the integral operator with eigenvalues $\mu_n = 1/n^2$. Watching $\lambda = n^2$ pass to $\mu = 1/n^2$ makes the operator-inverse relationship tactile: the unbounded spectrum of $L$ becomes the compact spectrum of $K$, accumulating at zero.

The same duality drives numerical practice. Finite-difference and finite-element methods discretize the differential form; the \emph{boundary element method} discretizes the integral form, trading a sparse matrix over a volume for a dense matrix over a boundary---exactly the local/global tradeoff made computational.

\section{In What Sense Are the Two Forms Equivalent?}
\label{sec:equivalence}

Mostly, but with caveats worth stating precisely.

\subsection{Regularity}

A solution of the integral equation need only be continuous (indeed, in more general settings, merely integrable), whereas a classical solution of the differential equation must possess enough derivatives. When $f$ is continuous, the integral equation bootstraps: a continuous solution is automatically $C^1$, hence (in the second-order case) $C^2$, and the two notions coincide. But when $f$ is merely measurable, the integral equation still makes sense while the pointwise differential equation does not. The integral formulation is the natural home of \emph{weak} and \emph{mild} solutions. This is not a defect; it is the standard way to \emph{define} solutions in low-regularity settings, and it is the gateway to the variational and weak-solution theory of partial differential equations.

\subsection{Direction of translation}

The passage from differential to integral form always works for a well-posed problem: integrate and absorb the data. The reverse passage works only when the kernel is smooth enough to differentiate through. A Volterra equation
\[
y(t) = g(t) + \int_{t_0}^{t} K(t, s)\, y(s)\, ds
\]
with smooth $K$ and $g$ can be differentiated back into a differential equation (or an integro-differential equation if $K$ depends nontrivially on $t$). But some integral equations have no local differential counterpart at all. Fredholm equations of the \emph{first kind},
\[
\int_a^b K(t, s)\, y(s)\, ds = g(t),
\]
in which the unknown appears only inside the integral, are the standard example: they ask one to invert a smoothing operator, an inherently unstable task, and they are genuinely ill-posed in Hadamard's sense.

\subsection{Locality versus globality}

This is the philosophical heart of the matter. A differential equation constrains behavior locally and requires side conditions to be supplied separately; the corresponding integral equation is a single global object containing law and data together. Neither is more fundamental. Physics moves freely between the two: the differential and integral forms of Maxwell's equations; the Schr\"odinger equation and the Lippmann--Schwinger equation of scattering theory.

\section{Summary}

The transformation from an initial value problem to a Volterra equation proceeds in four steps. First, integrate the differential equation from the data point $t_0$ to the generic point $t$; the limits are chosen so that the boundary term produced by the Fundamental Theorem of Calculus is a known quantity, and so that one obtains one equation per value of $t$. Second, rename the integration variable to avoid collision with the variable limit. Third, the Fundamental Theorem (Part 2) converts $\int y'$ into $y(t) - y_0$, absorbing the initial condition into the equation itself. Fourth, for the converse direction, the Fundamental Theorem (Part 1) differentiates the integral equation back into the differential equation, and setting $t = t_0$ recovers the data. Higher-order equations are handled by repeated integration---one auxiliary condition absorbed per integration---with Fubini's theorem collapsing the iterated integrals into a single integral against a polynomial kernel. Boundary value problems yield Fredholm equations whose kernels are Green's functions; and at the operator level the differential and integral formulations are inverse to one another, with the compact, smoothing integral side carrying the burden of rigorous analysis.

A good self-test is to run the whole procedure on $y' = y$, $y(0) = 1$, and then apply Picard iteration starting from $y_0(t) \equiv 1$: the partial sums of $e^t$ assemble themselves term by term, one integration per Taylor coefficient.

\end{document}
